\exercisesection{\S~1}

记 \(\mathfrak{s}\) 为李代数 \(\mathfrak{sl}(2,k)\)。

\begin{enumerate}
\def\labelenumi{\arabic{enumi})}
\tightlist
\item
  设 \(\mathrm{U}\) 为 \(\mathfrak{s}\) 的包络代数。证明元素
\end{enumerate}

\[
\mathrm{C} = H ^ {2} - 2 (X _ {+} X _ {-} + X _ {-} X _ {+}) = H ^ {2} + 2 H - 4 X _ {-} X _ {+}
\]

属于 U 的中心，并且它在与 \(V(m)\) 相关的表示中的像是比值为 \(m(m+2)\) 的伸缩变换。

\begin{enumerate}
\def\labelenumi{\arabic{enumi})}
\setcounter{enumi}{1}
\tightlist
\item
  设 \(\lambda \in k\)，设 \(Z(\lambda)\) 为具有基 \((e_n)\) 的向量空间，其中 \(n = 0,1,\ldots\)，并设 \(X_{+}, X_{-}, H\) 为由命题 1 的公式 (2) 定义的 \(Z(\lambda)\) 的自同态。
\end{enumerate}

\begin{enumerate}
\def\labelenumi{\alph{enumi})}
\item
  验证这给出一个\(\mathfrak{s}\)-模结构，作用于\(Z(\lambda)\)。
\item
  假设 \(\lambda\) 不是一个 \(\geq 0\) 的整数。证明 \(\mathfrak{s}\)-模 \(Z(\lambda)\) 是单的。
\item
  假设 \(\lambda\) 是一个 \(\geq 0\) 的整数。令 \(Z'\) 为 \(Z(\lambda)\) 的子空间，由 \(e_{n}\)（\(n > \lambda\)）生成。证明 \(Z'\) 是 \(Z(\lambda)\) 的 s-子模，同构于 \(Z(-\lambda - 2)\)，且商 \(Z(\lambda)/Z'\) 同构于单 s-模 \(V(\lambda)\)。证明 \(Z(\lambda)\) 的 s-子模只有 0、\(Z'\) 和 \(Z(\lambda)\)。
\end{enumerate}

\begin{enumerate}
\def\labelenumi{\arabic{enumi})}
\setcounter{enumi}{2}
\tightlist
\item
  设 \(\mathrm{E}\) 为具有基 \((e_n)_{n\in \mathbf{Z}}\) 的向量空间。令
\end{enumerate}

\[
a (n) = a _ {0} + a _ {1} n, b (n) = b _ {0} + b _ {1} n, c (n) = c _ {0} + c _ {1} n
\]

为系数属于 k 的三个仿射函数。按下列公式定义 E 的自同态 \(X_{+}, X_{-}\)、H：

\[
X _ {+} e _ {n} = a (n) e _ {n - 1}, X _ {-} e _ {n} = b (n) e _ {n + 1}, H e _ {n} = c (n) e _ {n}.
\]

当且仅当满足下列条件时，这给出 E 上的一个 \(\mathfrak{s}\)-模结构：

\[
a _ {1} b _ {1} = 1, c _ {1} = - 2, c _ {0} = - a _ {0} b _ {1} - a _ {1} b _ {0}.
\]

此 \(\mathfrak{s}\)-模在什么条件下是单的？

¶ 4) 设 g 为李代数。若 g-模 E 是有限维 g-子模的并，则称其为局部有限；等价地，E 的每个有限型（作为 U(g)-模的）g-子模都是有限维的。

\begin{enumerate}
\def\labelenumi{\alph{enumi})}
\item
  设 \(0 \rightarrow E' \rightarrow E \rightarrow E'' \rightarrow 0\) 是 g-模的正合序列。证明，如果 \(E'\) 和 \(E''\) 局部有限，则 E 也是局部有限（归约到 E 是有限型的情形，并利用 \(\mathrm{U}(\mathfrak{g})\) 是诺特环这一事实，参见~第 I 章，§2，节 6）。
\item
  假设 g 是半单的。证明 E 局部有限当且仅当 E 是有限维单 g-模的直和。
\item
  假设 \(\mathfrak{g} = \mathfrak{s}\)。证明 E 局部有限当且仅当满足以下条件：
\end{enumerate}

\(c_{1})\) E 有一个由 H 的特征向量组成的基，

\(c_{2}\) ) 自同态 \(X_{+\mathrm{E}}\) 和 \(X_{-\mathrm{E}}\) 局部幂零。

（先证明：如果 E 满足 \(c_{1}\) 和 \(c_{2}\)，且 E 不等于 0，则它包含一个具有整数权 \(m \geq 0\) 的本原元 e；再证明 \(X^{m+1}e = 0\)，从而 E 包含 \(V(m)\)。最后将 a) 应用于 E 的最大局部有限子模 \(E'\)。）

\begin{enumerate}
\def\labelenumi{\arabic{enumi})}
\setcounter{enumi}{4}
\tightlist
\item
  设 E 是一个局部有限的 \(\mathfrak{s}\)-模（习题 4）。对所有 \(m \geq 0\)，记 \(\mathrm{L}(m)\) 为 \(\mathfrak{s}\)-同态从 \(\mathrm{V}(m)\) 到 E 所成的向量空间。
\end{enumerate}

\begin{enumerate}
\def\labelenumi{\alph{enumi})}
\item
  定义从 E 到 \(\bigoplus_{m} \mathrm{L}_{m} \otimes \mathrm{V}(m)\) 的一个同构。
\item
  记 \(\Phi_{m}\) 为第 3 节备注 3 中定义的 \(\mathrm{V}(m)\) 上的不变双线性型。对所有 \(m \in N\)，令 \(b_{m}\) 为 \(L_{m}\) 上的双线性型；令 b 为 E 上的双线性型，它在 a) 中的同构下对应于双线性型 \(b_{m} \otimes \Phi_{m}\) 的直和。证明 b 是不变的，并且 E 上的每个不变双线性型都以唯一方式由此得到。b 是对称的（分别为交错的）当且仅当 m 为偶数时的 \(b_{m}\) 是对称的（分别为交错的），而 m 为奇数时的 \(b_{m}\) 是交错的（分别为对称的）。b 非退化当且仅当 \(b_{m}\) 非退化。
\item
  若 E 是有限维的，证明 E 是单生成的（作为 \(\mathrm{U}(\mathfrak{s})\)-模）当且仅当 \(\dim L_{m} \leq m + 1\) 对所有 \(m \geq 0\) 都成立。
\end{enumerate}

\begin{enumerate}
\def\labelenumi{\arabic{enumi})}
\setcounter{enumi}{5}
\tightlist
\item
  若 E 是有限维的 \(\mathfrak{s}\)-模，且 \(n\) 为一个整数，记 \(a_{n}\) 为 \(H_{\mathrm{E}}\) 的特征值为 \(n\) 的特征子空间的维数。记 \(c_{\mathrm{E}}(\mathrm{T})\) 为 \(\mathbf{Z}[\mathrm{T},\mathrm{T}^{-1}]\) 中由 \(c_{\mathrm{E}}(\mathrm{T}) = \sum_{n\in \mathbf{Z}}a_n\mathrm{T}^n\) 定义的元素。
\end{enumerate}

\begin{enumerate}
\def\labelenumi{\alph{enumi})}
\tightlist
\item
  按习题 5 定义 \(\mathrm{L}_m\)。证明
\end{enumerate}

\(\dim \mathrm{L}_m = a_m - a_{m + 2}\) 对每个整数 \(m\geq 0\)

推出，\(c_{\mathrm{E}}(\mathrm{T}) = c_{\mathrm{E}'}(\mathrm{T})\) 当且仅当 E 与 \(E'\) 同构。利用第 VII 章 §3 的习题 18 e) 重新得到这一结果。

\begin{enumerate}
\def\labelenumi{\alph{enumi})}
\setcounter{enumi}{1}
\item
  证明 \(c_{E\oplus F} = c_{E} + c_{F}\) 和 \(c_{E\otimes F} = c_{E}.c_{F}\)。
\item
  有 \(c_{\mathrm{V}(m)}(\mathrm{T}) = (\mathrm{T}^{m + 1} - \mathrm{T}^{-m - 1}) / (\mathrm{T} - \mathrm{T}^{-1})\)。
\end{enumerate}

\(d)\) 由 \(a\)、\(b\)、\(c)\) 推出：如果 \(m \geq m' \geq 0\)，则 \(\mathfrak{s}\)-模 \(\mathrm{V}(m) \otimes \mathrm{V}(m')\) 同构于

\[
\mathrm{V} (m + m ^ {\prime}) \oplus \mathrm{V} (m + m ^ {\prime} - 2) \oplus \mathrm{V} (m + m ^ {\prime} - 4) \oplus \dots \oplus \mathrm{V} (m - m ^ {\prime}).
\]

若 \(m\) 为 \(\geq 0\)，则 \(\mathfrak{s}\)-模 \(\mathbf{S}^2\mathrm{V}(m)\) 同构于

\[
\mathrm{V} (2 m) \oplus \mathrm{V} (2 m - 4) \oplus \mathrm{V} (2 m - 8) \oplus \dots \oplus \left\{ \begin{array}{l l} \mathrm{V} (0) & \text {if m is even} \\ \mathrm{V} (2) & \text {if m is odd} \end{array} \right.
\]

且 \(\mathfrak{s}\)-模 \(\bigwedge^2\mathrm{V}(m)\) 同构于

\[
\mathrm{V} (2 m - 2) \oplus \mathrm{V} (2 m - 6) \oplus \mathrm{V} (2 m - 10) \oplus \dots \oplus \left\{ \begin{array}{l l} \mathrm{V} (2) & \text {if m is even} \\ \mathrm{V} (0) & \text {if m is odd.} \end{array} \right.
\]

\begin{enumerate}
\def\labelenumi{\arabic{enumi})}
\setcounter{enumi}{6}
\item
  设 E 是有限维 s-模。证明 E 的对偶同构于 E。
\item
  证明 \(\mathfrak{s}\)-模 \(V(m)\) 可实现为两个变量中的齐次多项式 \(f(u,v)\)（次数为 \(m\)）所成的空间，其中算子 \(X_{+}, X_{-}\) 和 \(H\) 由下式给出：
\end{enumerate}

\[
X _ {+} f = u \frac {\partial f}{\partial v}, X _ {-} f = - v \frac {\partial f}{\partial u}, H f = u \frac {\partial f}{\partial u} - v \frac {\partial f}{\partial v}.
\]

¶9) 考虑 \(\mathfrak{s}\) 通过伴随表示在其包络代数 U 和其对称代数 \(\mathbf{S} = \bigoplus \mathbf{S}^n\) 上的作用。

\begin{enumerate}
\def\labelenumi{\alph{enumi})}
\tightlist
\item
  确定 \(\mathfrak{s}\)-模 \(\mathbf{S}^n\) 的权。推出以下 \(\mathfrak{s}\)-模同构：
\end{enumerate}

\[
\begin{array}{l l} \mathbf {S} ^ {n} \longrightarrow \mathrm{V} (2 n) \oplus \mathrm{V} (2 n - 4) \oplus \mathrm{V} (2 n - 8) \oplus \dots \oplus \mathrm{V} (0) & n \text {even} \\ \mathbf {S} ^ {n} \longrightarrow \mathrm{V} (2 n) \oplus \mathrm{V} (2 n - 4) \oplus \mathrm{V} (2 n - 8) \oplus \dots \oplus \mathrm{V} (2) & n \text {odd}. \end{array}
\]

特别地，\(\mathbf{S}^n\) 中在 \(\mathfrak{s}\) 下不变的元素构成一个空间；当 \(n\) 为偶数（分别为奇数）时，其维数为 1（分别为 0）。

\begin{enumerate}
\def\labelenumi{\alph{enumi})}
\setcounter{enumi}{1}
\item
  证明由在 s 下不变的元素组成的 S 的子代数是多项式代数 \(k[\Gamma]\)，其中 \(\Gamma = H^{2} - 4X_{-}X_{+}\)。模 \(S^{n}/\Gamma.S^{n-2}\) 同构于 \(\mathrm{V}(2n)\)。
\item
  证明 U 的中心是多项式代数 \(k[\mathbf{C}]\)，其中 C 是习题 1 中定义的元素（利用 b) 和 Poincaré-Birkhoff-Witt 定理）。
\end{enumerate}

\begin{enumerate}
\def\labelenumi{\arabic{enumi})}
\setcounter{enumi}{9}
\tightlist
\item
  设 \(m\) 为一个 \(> 0\) 的整数，\(\mathbf{S}\) 为分次代数 \(k[\mathrm{X}_1, \ldots, \mathrm{X}_m]\)，且 \(a_1, \ldots, a_m\) 为 \(k^*\) 中的元素。设 \(\Phi = \sum_{i,j=1}^{m} a_{ij} X_i X_j\) 为一个二次型；置 \(\mathrm{D}_i = \frac{\partial}{\partial X_i}\)，并置 \(\mathrm{D} = \sum_{i,j=1}^{m} b_{ij} \mathrm{D}_i \mathrm{D}_j\)，其中 \((b_{ij})\) 是 \((a_{ij})\) 的逆矩阵。
\end{enumerate}

\begin{enumerate}
\def\labelenumi{\alph{enumi})}
\tightlist
\item
  证明存在唯一的 \(\mathfrak{s}\)-模结构，作用在 \(\mathbf{S}\) 上，使得 \(X_{+}f = \frac{1}{2}\mathrm{D}(f), X_{-}f = \frac{1}{2}\Phi f\)，并且 \(Hf = -(m / 2 + n)f\)（当 \(f\) 是次数为 \(n\) 的齐次元素时）。b) 置 \(\mathrm{A}^n = \mathbf{S}^n \cap \mathrm{Ker}(\mathrm{D})\)。证明 \(\mathbf{S}^n\) 是 \(\Phi^p\)、\(\mathrm{A}^q\)（其中 \(2p + q = n\)）这些子空间的直和。推出恒等式
\end{enumerate}

\[
\sum_ {n = 0} ^ {\infty} \dim (\mathrm{A} ^ {n}) \mathrm{T} ^ {n} = (1 + \mathrm{T}) / (1 - \mathrm{T}) ^ {m - 1}.
\]

\begin{enumerate}
\def\labelenumi{\alph{enumi})}
\setcounter{enumi}{2}
\item
  若 \(f\) 是 \(\mathbf{A}^n\) 的非零元素，则 \(\varPhi^p f\)（\(p \geq 0\)）构成一个单的 \(\mathfrak{s}\)-子模的基；该子模属于 \(\mathbf{S}\)，并同构于习题 2 中的模 \(Z\left(-\frac{m}{2} - n\right)\)。
\item
  具体写出 \(m = 1\) 和 \(m = 2\) 的情形。利用 \(m = 3\) 的情形重新得到习题 9 的结果。
\end{enumerate}

\begin{enumerate}
\def\labelenumi{\arabic{enumi})}
\setcounter{enumi}{10}
\tightlist
\item
  假设 \(k\) 是代数闭的。设 \(\mathfrak{g}\) 为李代数，\(V\) 为单 \(\mathfrak{g}\)-模，且 \(D\) 为 \(\mathfrak{g}\)-模 \(V\) 的自同态域。证明，如果 \(k\) 是不可数的，\(^7\) 则 \(D = k\)。（否则，\(D\) 将包含一个同构于 \(k(X)\) 的子域，其中 \(X\) 是不定元，于是有 \(\dim_k D > \aleph_0\)，从而也有 \(\dim_k V > \aleph_0\)，这与 \(V\) 是单生成的 \(U(\mathfrak{g})\)-模相矛盾。）
\end{enumerate}

\footnotetext[7]{结论即使 k 可数时仍然成立，参见 D. QUILLEN，《关于包络代数上一个单模的自同态环》，Proc. Amer. Math. Soc., Vol. XXI (1969), pp. 171-172。}

¶ 12) 设 \(q \in k\)，且 \(W\) 为具有基 \((e_0, e_1, e_2, \ldots)\) 的向量空间。

\begin{enumerate}
\def\labelenumi{\alph{enumi})}
\tightlist
\item
  证明存在唯一的表示 \(\rho_{q}\)，它是 \(\mathfrak{s}\) 上的表示并作用在 W 上，使得
\end{enumerate}

\[
\begin{array}{r l} & {\rho_ {q} (H) e _ {n} = 2 e _ {n + 1}, \quad \rho_ {q} (X _ {+}) e _ {n} = (\frac {1}{2} \rho_ {q} (H) - 1) ^ {n} e _ {0},} \\ & {\rho_ {q} (X _ {-}) e _ {n} = (\frac {1}{2} \rho_ {q} (H) + 1) ^ {n} (- q e _ {0} + e _ {1} + e _ {2}).} \end{array}
\]

有 \(\rho_{q}(\mathrm{C}) = 4q\)，其中 \(C = H^{2} + 2H - 4X_{-}X_{+}\)（参见~习题 1）。表示 \(\rho_{q}\) 是单的。元素 \(x \in s\) 使 \(\rho_{q}(x)\) 具有特征值的，正是 \(X_{+}\) 的倍数。自同态 \(\rho_{q}(X_{+})\) 以重数 1 具有特征值 1。

\begin{enumerate}
\def\labelenumi{\alph{enumi})}
\setcounter{enumi}{1}
\tightlist
\item
  设 \(\rho\) 是 \(\mathfrak{s}\) 的一个单表示，满足 \(\rho(\mathrm{C}) = 4q\)，且 \(\rho(X_{+})\) 具有特征值 1。证明 \(\rho\) 等价于 \(\rho_q\)。
\end{enumerate}

¶ 13) 假设 \(k = \mathbf{C}\)。置 \(C = H^2 + 2H - 4X_-X_+\)，参见~习题 1。一个表示 \(\rho\)（其对象为 \(\mathfrak{s} = \mathfrak{sl}(2, \mathbf{C})\)）被称为 \(H\)-可对角化的，如果 \(\rho\) 的底层空间有一个由 \(\rho(H)\) 的特征向量组成的基。令 \(q \in \mathbf{C}\)，\(v \in \mathbf{Q}/\mathbf{Z}\)。

\begin{enumerate}
\def\labelenumi{\alph{enumi})}
\tightlist
\item
  设 S 为具有基 \((e_{w})_{w\in\mathbf{C}}\) 的向量空间，该基由 C 的元素索引。令 \(S_{v}=\sum_{w\in v}Ce_{w}\)。存在唯一的表示 \(\rho_{v,q}\)，它是 s 在 \(S_{v}\) 上的表示，使得
\end{enumerate}

\[
\rho_ {v, q} (X _ {+}) e _ {w} = (q - w ^ {2} - w) ^ {1 / 2} e _ {w + 1}
\]

\[
\begin{array}{c} \rho_ {v, q} (X _ {-}) e _ {w} = (q - w ^ {2} + w) ^ {1 / 2} e _ {w - 1} \\ \rho_ {v, q} (H) e _ {w} = 2 w e _ {w}. \end{array}
\]

（约定，对于所有 \(z \in \mathbf{C}\)，\(z^{1/2}\) 是 \(z\) 的平方根，其幅角属于 \(\{0, \pi\}\)。）记 \(\mathrm{S}_{v,q}\) 为一个 \(\mathfrak{s}\)-模，它由 \(\rho_{v,q}\) 定义。有 \(\rho_{v,q}(\mathrm{C}) = 4q\)。

\begin{enumerate}
\def\labelenumi{\alph{enumi})}
\setcounter{enumi}{1}
\item
  若 \(q \neq u^2 + u\) 对所有 \(u \in v\) 都成立，则 \(S_{v,q}\) 是单的。
\item
  假设 \(2v \neq 0\)，且 q 形如 \(u^{2} + u\)，其中 \(u \in v, u \geq 0\)（这唯一确定 u）。令 \(S_{v,q}^{-}\)（分别为 \(S_{v,q}^{+}\)）为 \(S_{v}\) 中由 \(e_{w}\)（满足 \(w \leq u\)；分别为 w \textgreater{} u）生成的向量子空间。则 \(S_{v,q}^{-}\) 和 \(S_{v,q}^{+}\) 是 \(S_{v,q}\) 的单 s-子模。
\item
  假设 2v = 0，且 q 形如 \(u^{2} + u\)，其中 \(u \in v, u \geq 0\)（这唯一确定 u）。令 \(S_{v,q}^{-}\)（分别为 \(S_{v,q}^{0}, S_{v,q}^{+}\)）为 \(S_{v}\) 中由 \(e_{w}\)（满足 w \textless{} -u；分别为 \(-u \leq w \leq u, w > u\)）生成的向量子空间。则 \(S_{v,q}^{-}, S_{v,q}^{0}\) 和 \(S_{v,q}^{+}\) 是 \(S_{v,q}\) 的单 s-子模。
\item
  假设 \(v = -\frac{1}{2} + Z\) 且 \(q = -\frac{1}{4}\)。令 \(S_{-1/2,-1/4}^{-}\)（分别为 \(S_{-1/2,-1/4}^{+}\)）为 \(S_{-1/2}\) 中由 \(e_{w}\)（满足 \(w \leq -\frac{1}{2}\)；分别为 \(w > -\frac{1}{2}\)）生成的向量子空间。则 \(S_{-1/2,-1/4}^{-}\) 和 \(S_{-1/2,-1/4}^{+}\) 是 \(S_{-1/2,-1/4}\) 的单 s-子模。
\end{enumerate}

\(f)\) 记 \(\rho_{v,q}^{\pm},\rho_{v,q}^{0}\) 为分别与 \(\mathrm{S}_{v,q}^{\pm},\mathrm{S}_{v,q}^{0}\) 对应的表示。在情形 \(b)\) 中，元素 \(x\in \mathfrak{s}\) 若使 \(\rho_{v,q}^{-}(x)\) 具有特征值，则这些元素正是 \(\mathbf{CH}\) 中的元素。在情形 \(c)\)、\(d)\)、\(e)\) 中，元素 \(x\in \mathfrak{s}\) 若使 \(\rho_{v,q}^{-}(x)\) 具有特征值，则这些元素正是 \(\mathbf{CH} + \mathbf{CX}_+\) 中的元素；此外，若 \(x\) 是幂零的（因而与 \(X_{+}\) 成比例）且非零，则 \(\rho_{v,q}^{-}(x)\) 的唯一特征值为 0，且其重数为 1；另一方面，若 \(x\) 是半单的，则 \(\rho_{v,q}^{-}\) 的底层空间有一个由 \(\rho_{v,q}^{-}(x)\) 的特征向量组成的基。对于 \(\rho_{v,q}^{+}\) 有类似结果，只需将 \(X_{+}\) 换成 \(X_{-}\)。

\(g)\) 设 \(V\) 是单的 \(\mathfrak{s}\)-模，\(\rho\) 为相应表示。则 \(\rho(C)\) 是伸缩变换（利用习题 11）。假设 \(\rho(C) = 4q\)。证明，如果 \(\rho\) 是 \(H\)-可对角化的，则 \(V\) 同构于模 \(S_{v,q}, S_{v,q}^{\pm}, S_{v,q}^{0}\)，该模在 \(b), c), d), e)\) 中被考虑。此外，\(\rho\) 是 \(H\)-可对角化的当且仅当 \(\rho(H)\) 具有特征值；只要 \(\rho(X_{+})\) 具有特征值 \(0\) 即可。

¶14) 记号沿用上一个习题。记 \(B_{q}\) 为 \(U(\mathfrak{s})\) 关于由 C - 4q 生成的双边理想的商，并记 \(u \mapsto u^{\bullet}\) 为 \(U(\mathfrak{s}) \to B_{q}\) 的典则映射。将习题 12 和 13 的表示看作 \(B_{q}\) 的表示。

\begin{enumerate}
\def\labelenumi{\alph{enumi})}
\tightlist
\item
  \(B_{q}\) 的每个元素都能唯一表示为
\end{enumerate}

\[
\sum_ {r \geq 0} X _ {+} ^ {\bullet r} p _ {r} (H ^ {\bullet}) + \sum_ {s > 0} q _ {s} (H ^ {\bullet}) X _ {-} ^ {\bullet s},
\]

其中 \(p_{r}\) 和 \(q_{s}\) 是多项式。若 \(B_{q}\) 的两个元素生成同一个左理想，则它们成比例。

\begin{enumerate}
\def\labelenumi{\alph{enumi})}
\setcounter{enumi}{1}
\item
  考虑习题 13 的情形 b)。设 a 是 \(S_{v,q}\) 的非零元素。如果 a 是 \(\rho_{v,q}(H)\) 的特征向量，特征值为 \(\lambda\)，则 a 在 \(B_{q}\) 中的湮没子是由 \(H^{\bullet}-\lambda\) 生成的左理想；如果 a 不是 \(\rho_{v,q}(H)\) 的特征向量，则其湮没子是非单生成左理想。
\item
  考虑习题 13 的表示 \(\rho_{v,q}^{\pm}\) 的情形。如果 a 是这种表示的空间中的非零元素，则其在 \(B_{q}\) 中的湮没子是非单生成左理想。
\item
  考虑习题 12 的表示 \(\rho_{q}\)。\(e_{0}\) 在 \(B_{q}\) 中的湮没子由 \(X_{+}^{\bullet}-1\) 生成。如果 \(a\in W\) 不与 \(e_{0}\) 成比例，则其湮没子是非单生成左理想。
\item
  每个 \(\mathfrak{s}\) 的自同构都延拓为 \(\mathrm{U}(\mathfrak{s})\) 的自同构，并通过取商定义 \(\mathrm{B}_q\) 的一个自同构；令 \(\mathrm{A}'\) 为由此得到的 \(\mathrm{A} = \mathrm{Aut}(\mathrm{B}_q)\) 的子群。证明 \(\mathrm{A}' \neq \mathrm{A}\)。（令 \(\varphi\) 为端同态 \(x \mapsto [X_{+}^{\bullet 2}, x]\)；它是 \(\mathrm{B}_q\) 的局部幂零自同态，且 \(e^{\varphi} \in \mathrm{A}, e^{\varphi} \notin \mathrm{A}'\)。）
\item
  群 A 通过结构迁移作用在 \(B_{q}\) 的单表示类集合上。令 \(\pi_{1}\)（分别为 \(\pi_{2},\pi_{3}\)）为习题 13 b) 的 \(\rho_{v,q}\) 型表示（分别为习题 13 的 \(\rho_{v,q}^{\pm}\) 型表示，分别为习题 12 的 \(\rho_{q}\) 型表示）。则 \(A\pi_{1},A\pi_{2}\) 和 \(A\pi_{3}\) 两两不相交。（利用 a)、b)、c)、d)。若 \(\psi\in A\) 使得 \(\psi\pi_{1}\) 是习题 13 b) 的 \(\rho_{v,q}\) 型，则 \(\psi\in A'\)（利用 a) 和 b)）。推出，若 \(\psi\) 是 e) 中的自同构 \(e^{\varphi}\)，则表示 \(\sigma=\pi_{1}\circ\psi\) 具有如下性质：对所有 \(x\in s-\{0\}\)，\(\sigma(x)\) 都没有特征值。 \(^{8}\)
\end{enumerate}

\footnotetext[8]{关于习题 12、13 和 14 的更多细节，参见：D. ARNAL 和 G. PINCZON，《关于 sl(2) 李代数的代数不可约表示》，J. Math. Phys., Vol. 15 (1974), pp. 350-359。}

\begin{enumerate}
\def\labelenumi{\arabic{enumi})}
\setcounter{enumi}{14}
\tightlist
\item
  设 \(\mathfrak{g}\) 是一个 3 维李代数。证明以下条件等价（参见~第 I 章，§6，习题 23）。
\end{enumerate}

\begin{enumerate}
\def\labelenumi{(\roman{enumi})}
\item
  \(\mathfrak{g} = [\mathfrak{g},\mathfrak{g}]\)。
\item
  \(\mathfrak{g}\) 的 Killing 型非退化。
\item
  g 是半单的。
\item
  g 是单的。
\end{enumerate}

¶ 16) 设 g 是一个 3 维单李代数，且 \(\Phi\) 是它的 Killing 型。记 \(\mathfrak{o}(\Phi)\) 为 \(\Phi\) 的正交代数，即由 \(\mathfrak{gl}(\mathfrak{g})\) 中保持 \(\Phi\) 不变的元素组成的子代数。

\begin{enumerate}
\def\labelenumi{\alph{enumi})}
\item
  证明 ad：\(\mathfrak{g} \to \mathfrak{o}(\varPhi)\) 是一个同构。
\item
  证明以下性质等价：
\end{enumerate}

\begin{enumerate}
\def\labelenumi{(\roman{enumi})}
\item
  \(\mathfrak{g}\) 含有一个非零迷向向量（相对于 \(\varPhi\)）。
\item
  \(\mathfrak{g}\) 含有一个非零幂零元。
\item
  g 同构于 s。
\end{enumerate}

\begin{enumerate}
\def\labelenumi{\alph{enumi})}
\setcounter{enumi}{2}
\item
  证明存在一个扩域 \(k_{1}\)，其基域为 \(k\)，次数为 \(\leq 2\)，使得 \(\mathfrak{g}_{(k_1)} = k_1 \otimes_k \mathfrak{g}\) 同构于 \(\mathfrak{s}_{(k_1)}\)。
\item
  在向量空间 \(\mathbf{A} = k\oplus \mathfrak{g}\) 上赋予唯一的代数结构，使 1 为单位元，并使得乘法 \(\mathfrak{g}\times \mathfrak{g}\to \mathrm{A}\) 由公式给出
\end{enumerate}

\[
x. y = \frac {1}{8} \Phi (x, y). 1 + \frac {1}{2} [ x, y ] \quad (x, y \in \mathfrak {g}).
\]

证明 A 是 k 上的四元数代数，而 g 是 A 中由迹为零的元素组成的李子代数（通过扩张基域归约到 g = s 的情形，并证明此时 A 可与矩阵代数 \(\mathbf{M}_{2}(k)\) 认同）。

反之，若 D 是 k 上的四元数代数，则 D 中迹为零的元素构成一个 3 维单李代数，并且相应的代数 A 可与 D 认同。

\begin{enumerate}
\def\labelenumi{\alph{enumi})}
\setcounter{enumi}{4}
\tightlist
\item
  证明下列公式
\end{enumerate}

\[
\begin{array}{c} \Phi (x, x) = - 8 \mathrm{Nrd} _ {\mathrm{A}} (x), \quad \Phi (x, y) = 4 \mathrm{Trd} _ {\mathrm{A}} (x y) \\ 2 [ x, [ y, z ] ] = \Phi (x, y) z - \Phi (x, z) y \qquad (x, y, z \in \mathfrak {g}). \end{array}
\]

\(f\) ) 证明 \(\varPhi\) 的判别式（相对于 \(\mathfrak{g}\) 的任意基）具有 \(-2\lambda^2\) 的形式，其中 \(\lambda \in k^{*}\)。

\begin{enumerate}
\def\labelenumi{\alph{enumi})}
\setcounter{enumi}{6}
\tightlist
\item
  设 n 是一个 \(\geq 0\) 的整数。证明 g-模 \(\mathbf{S}^{n}(\mathfrak{g})\) 有一个唯一的维数为 \(2n + 1\) 的单子模，且此模是绝对单的（归约到 g = s 的情形，并利用习题 9）。
\end{enumerate}

证明 g 只有在同构于 s（即~只有在 A 同构于 \(\mathbf{M}_{2}(k)\)）时才有维数为 2n 的绝对单模。（设 V 是这样的模。如果 \(n \geq 2\)，利用习题 6 c) 证明 g-模 \(V \otimes g\) 有唯一的维数为 2n - 2 的绝对单子模。然后归约到 n = 1 的情形，这显然。）

¶ 17) 保留上一个习题的记号。证明，对所有 \(n \geq 1\)，代数 \(\mathrm{U}(\mathfrak{g})\) 有唯一的双边理想 \(m_{n}\)，使得 \(\mathrm{U}(\mathfrak{g}) / \mathfrak{m}_{n}\) 是维数为 \(n^{2}\) 的中心单代数（扩张标量以归约到 g = s 的情形，并证明此时 \(m_{n}\) 是同态 \(\mathrm{U}(\mathfrak{g}) \to \mathrm{End}(\mathrm{V}(n-1))\) 的核）。\(\mathrm{U}(\mathfrak{g})\) 的每个有限余维双边理想都形如 \(m_{n_{1}} \cap m_{n_{2}} \cap \cdots \cap m_{n_{h}}\)，其中 \(n_{1}, \ldots, n_{h}\) 互异；其余维为 \(n_{1}^{2} + \cdots + n_{h}^{2}\)（应用稠密性定理）。\(m_{n}\) 是 \(\mathrm{U}(\mathfrak{g})\) 中唯一的有限余维极大双边理想。

证明 \(m_{2}\)（作为双边理想）由元素 \(x^{2}-\frac{1}{8}\Phi(x,x)\)（\(x\in g\)）生成，并且商 \(\mathrm{U}(\mathfrak{g})/\mathfrak{m}_{2}\) 可与习题 16 的四元数代数 A 认同。

当 g 同构于 s 时，\(\mathrm{U}(\mathfrak{g})/\mathfrak{m}_{n}\) 同构于 \(\mathbf{M}_{n}(k)\)。利用习题 16 g)，证明当 g 不同构于 s 时，\(\mathrm{U}(\mathfrak{g})/\mathfrak{m}_{n}\) 同构于 \(\mathbf{M}_{n}(k)\) 当且仅当 n 为奇数。\(^{9}\)

\footnotetext[9]{可以证明，当 n 为偶数时，\(\mathrm{U}(\mathfrak{g})/\mathfrak{m}_{n}\) 同构于 \(\mathbf{M}_{n/2}(\mathrm{A})\)。}

¶ 18) 假设 \(k = \mathbf{R}, \mathbf{C}\)，或者 k 是一个关于离散赋值完备且剩余域特征为 \(p \neq 0\) 的域（例如 \(p\)-进域 \(\mathbf{Q}_p\) 的有限扩域）。

\begin{enumerate}
\def\labelenumi{\alph{enumi})}
\tightlist
\item
  设 \(\mathfrak{n}\) 是一个幂零李代数，N 是通过赋予 \(\mathfrak{n}\) Hausdorff 作用律得到的李群（第 III 章，§9，第 5 节），且 \(\rho : \mathfrak{n} \to \mathfrak{gl}(\mathrm{E})\) 是 \(\mathfrak{n}\) 在有限维向量空间 E 上的线性表示。假设 \(\rho(x)\) 对所有 \(x \in \mathfrak{n}\) 都是幂零的，并置 \(\pi(x) = \exp(\rho(x))\)。证明 \(\pi\) 是唯一的李群同态 \(\varphi : \mathrm{N} \to \mathbf{GL}(\mathrm{E})\)，使得 \(\mathrm{L}(\varphi) = \rho\)。
\end{enumerate}

（当 k = R 或 C 时，利用 N 连通这一事实。当 k 是超度量的时，证明 \(\varphi(x)\)（\(x \in N\)）的特征值等于 1；为此，先证明：如果 \(k'\) 是 k 的有限扩域，且 \((\lambda_{1}, \ldots, \lambda_{n}, \ldots)\) 是 \(k'\) 中元素的序列，满足对所有 n 有 \(\lambda_{n} = \lambda_{n+1}^{p}\)，且 \(\lambda_{1} = 1\)，则对所有 n 有 \(\lambda_{n} = 1\)。）

\begin{enumerate}
\def\labelenumi{\alph{enumi})}
\setcounter{enumi}{1}
\tightlist
\item
  设 \(\rho : \mathfrak{s} \to \mathfrak{gl}(\mathrm{E})\) 是 \(\mathfrak{s}\) 的有限维线性表示，并设 \(\pi\) 是从 \(\mathbf{SL}(2, k)\) 到 \(\mathbf{GL}(\mathrm{E})\) 的与 \(\rho\) 相容的同态（第 4 节）。证明 \(\pi\) 是唯一的李群同态 \(\varphi : \mathbf{SL}(2, k) \to \mathbf{GL}(\mathrm{E})\)，使得 \(\mathrm{L}(\varphi) = \rho\)。（利用 \(a\) ) 证明 \(\pi\) 与 \(\rho\) 在 \(\exp(n)\) 上一致，其中 \(n\) 是 \(\mathfrak{s}\) 中的幂零元，并注意到 \(\mathbf{SL}(2, k)\) 由 \(\exp(n)\) 生成。）
\end{enumerate}

